\begin{textsample}{POS Dim 4 – persona\_gpt – Score 29.00 – t637\_gpt.txt}  \label{ex:f4_pos_023}
Around the world ’s \textbf{oceans}, seabirds are sounding a clear warning about the damage humans are inflicting on marine \textbf{ecosystems}.
Among the most threatened are albatrosses, iconic \textbf{ocean} wanderers that can travel \textbf{thousands} of \textbf{kilometres} across the \textbf{seas}.
Their survival is being undermined on multiple fronts, with human \textbf{activities} reshaping the \textbf{oceans} faster than these long‑lived \textbf{birds} can adapt.
One of the most severe pressures is industrial overfishing.
By removing vast quantities of \textbf{fish} from the \textbf{sea}, industrial \textbf{fleets} are stripping away the prey that seabirds depend on.
For albatrosses and many other \textbf{species}, this means longer, riskier foraging trips and less food for their chicks.
Since 1950, some populations have \textbf{declined} by up to 70 %, a stark indicator of how quickly things can unravel when the food \textbf{web} is destabilised.
On top of food scarcity, many seabirds are being killed directly in \textbf{fishing} \textbf{operations}.
Bycatch – the unintentional capture of non‑target \textbf{species} – remains a major \textbf{threat}. \textbf{Birds} can become fatally entangled in lines and \textbf{nets} or drown when they seize baited hooks and are dragged underwater.
These deaths are entirely avoidable, yet they continue at scale in many \textbf{fisheries}.
Plastic pollution adds another deadly \textbf{layer} of pressure.
More than a million seabirds are killed each year by plastics, whether through ingestion or entanglement. \textbf{Birds} mistake floating fragments for food, filling their stomachs with material that provides no nutrition and can cause internal injuries or starvation.
Others become trapped in discarded \textbf{fishing} gear and other debris.
The sheer volume of plastic in the \textbf{oceans} is turning once‑rich feeding grounds into hazardous traps.
Seabirds are also vulnerable to \textbf{threats} on land and at \textbf{sea} that have built up over generations.
Invasive \textbf{species} introduced to nesting \textbf{islands} – such as rats, cats, and other predators – can devastate colonies by eating eggs, chicks, and even adults.
Historically, feather harvesting drove intense persecution, with \textbf{birds} killed in huge numbers for fashion and trade.
While that particular industry has waned, its legacy lingers in depleted populations and lost breeding strongholds.
Modern industrial society has introduced new dangers as well.
Seabirds collide with \textbf{ships} and structures at \textbf{sea}, and oil \textbf{spills} can be catastrophic, coating feathers, destroying insulation, and poisoning \textbf{birds} that ingest contaminated \textbf{water} or prey.
Toxic chemicals that accumulate in marine food \textbf{webs} pose additional, often invisible, risks to their health and reproductive success. \textbf{Layered} over all of this is the growing influence of climate change and \textbf{ocean} acidification.
As humans burn fossil fuels, carbon \textbf{dioxide} emissions are not only warming the planet but also changing the chemistry of the \textbf{oceans}.
Since the Industrial Revolution, \textbf{ocean} acidity has increased by about 30 %.
This shift makes \textbf{life} harder for organisms that build \textbf{shells} and skeletons, such as corals and shellfish.
When these foundational \textbf{species} are harmed, the effects ripple through entire \textbf{ecosystems}, altering \textbf{habitats} and food availability for countless marine \textbf{animals}, including seabirds.
Studies from around the world are \textbf{documenting} how these overlapping pressures are disrupting marine \textbf{ecosystems}.
Seabirds, which are highly mobile and closely tied to \textbf{ocean} productivity, are among the first to show the consequences.
Their \textbf{declines} are a visible sign that the systems they depend on – and that humans rely on for food, climate regulation, and \textbf{livelihoods} – are under strain.
Yet there are clear pathways to help seabird populations recover and to restore healthier \textbf{oceans}.
Expanding marine \textbf{reserves} can provide safe havens where industrial \textbf{fishing} is limited or banned, allowing prey \textbf{fish} populations to rebuild and reducing the risk of bycatch.
Banning destructive practices such as \textbf{bottom} trawling in sensitive \textbf{areas} can protect crucial feeding and breeding \textbf{habitats} from being literally scraped off the seafloor.
Restoration efforts, including the removal of invasive \textbf{species} from nesting \textbf{islands} and the rehabilitation of damaged \textbf{habitats}, can give seabirds a chance to rebound.
These actions are not just about saving individual \textbf{species} ; they are about rebuilding the resilience of entire marine \textbf{ecosystems}.
Seabirds are powerful indicators of \textbf{ocean} health.
When they thrive, it is a sign that the \textbf{seas} are rich, balanced, and capable of sustaining \textbf{life}.
Protecting them means confronting overfishing, plastic pollution, fossil fuel dependence, and the many other pressures we have placed on the \textbf{oceans} – and choosing instead to give marine \textbf{life} the space and time it needs to recover.

% matched lemmas: activity, animal, area, bird, bottom, decline, dioxide, document, ecosystem, fish, fishery, fishing, fleet, habitat, island, kilometre, layer, life, livelihood, net, ocean, operation, reserve, sea, shell, ship, specie, spill, thousand, threat, water, web
\end{textsample}
